\documentclass[11pt,a4paper]{article}

\usepackage[utf8]{inputenc}
\usepackage[T1]{fontenc}
\usepackage[spanish,es-noshorthands]{babel}
\usepackage{xcolor}
\usepackage[most]{tcolorbox}
\usepackage{listings}
\usepackage{fancyhdr}
\usepackage{lastpage}
\usepackage{enumitem}
\usepackage[hidelinks]{hyperref}

\setlength{\headheight}{30pt}
\pagestyle{fancy}
\fancyhf{}
\renewcommand{\headrulewidth}{0pt}
\fancyhead[R]{%
  \fontsize{8}{9.6}\selectfont
  Licenciatura en Ciencias de la Computación\\
  Sistemas Operativos\\[2pt]
  \rule{4.2cm}{0.5pt}
}
\fancyfoot[C]{\thepage\ / \pageref{LastPage}}

\newtcolorbox{infobox}{
  colback=blue!6!white, colframe=blue!45!white,
  arc=2mm, boxrule=0.6pt, left=8pt, right=8pt, top=6pt, bottom=6pt,
  title=INFO, fonttitle=\bfseries, coltitle=blue!45!black,
  colbacktitle=blue!6!white, boxsep=1pt, breakable
}

\newtcolorbox{extrabox}{
  colback=orange!8!white, colframe=orange!55!white,
  arc=2mm, boxrule=0.6pt, left=8pt, right=8pt, top=6pt, bottom=6pt,
  title=EXTRA, fonttitle=\bfseries, coltitle=orange!55!black,
  colbacktitle=orange!8!white, boxsep=1pt, breakable
}

\newtcolorbox{notebox}{
  colback=black!5!white, colframe=black!25!white,
  arc=2mm, boxrule=0.5pt, left=8pt, right=8pt, top=6pt, bottom=6pt,
  boxsep=1pt, breakable
}

\lstdefinestyle{codestyle}{
  backgroundcolor=\color{black!5},
  rulecolor=\color{black!25},
  frame=single,
  framerule=0.5pt,
  basicstyle=\ttfamily\small,
  breaklines=true,
  showstringspaces=false,
  tabsize=4,
  keywordstyle=\color{blue!60!black}\bfseries,
  commentstyle=\color{green!40!black},
  stringstyle=\color{orange!70!black},
  columns=fullflexible,
  upquote=true
}
\lstset{style=codestyle}
\newcommand{\inputcode}[2][]{%
  \par\addvspace{0.2\baselineskip}\noindent
  \begin{minipage}{\linewidth}
    \lstinputlisting[#1]{#2}
  \end{minipage}\par\addvspace{0.2\baselineskip}}
\setlist{nosep}

\title{Trabajo Práctico 5 -- Sincronización}
\author{}
\date{}

\begin{document}
\maketitle
\thispagestyle{fancy}
\begin{center}
SISTEMAS OPERATIVOS
\end{center}
\newpage

\tableofcontents
\newpage

En este TP vas a estudiar cómo cooperan los procesos cuando deben
comunicarse o coordinar el acceso a un recurso compartido. Vamos a usar
llamadas POSIX (\emph{Portable Operating System Interface}) para construir una
tubería y candados de archivos, y la interfaz de IPC de System V para los
semáforos. En todos los casos, observá primero qué hace el núcleo y después
relacioná el resultado con los descriptores, los procesos y sus estados.

\section{Comunicación entre procesos con tuberías}

Una tubería o \emph{pipe} es un canal unidireccional de comunicación entre
procesos. Un proceso escribe una secuencia de bytes en un extremo y otro
proceso la lee del extremo opuesto. Por ejemplo, el shell puede conectar
\texttt{ls} con \texttt{wc} mediante el carácter \texttt{|}:

\begin{lstlisting}[language=bash]
$ ls | wc
\end{lstlisting}

\texttt{wc} (\emph{word count}) muestra, por defecto, la cantidad de líneas,
palabras y bytes que recibió. El shell no copia esos datos a mano: crea un
pipe, conecta los descriptores estándar de cada proceso y luego ejecuta los
programas.

\subsection{Descriptores de archivo}

Un descriptor de archivo es un número entero que identifica, dentro de un
proceso, un archivo o canal abierto. Cada proceso comienza normalmente con
tres descriptores estándar:

\begin{center}
\begin{tabular}{c|l|l}
\textbf{Valor} & \textbf{Nombre} & \textbf{Constante en C} \\
\hline
0 & Entrada estándar (\texttt{stdin}) & \texttt{STDIN\_FILENO} \\
1 & Salida estándar (\texttt{stdout}) & \texttt{STDOUT\_FILENO} \\
2 & Error estándar (\texttt{stderr}) & \texttt{STDERR\_FILENO}
\end{tabular}
\end{center}

La llamada \texttt{pipe(fd)} recibe un arreglo de dos enteros. Si tiene
éxito, \texttt{fd[0]} es el extremo de lectura y \texttt{fd[1]} el extremo de
escritura. Un pipe anónimo no tiene un nombre en el sistema de archivos: los
procesos pueden usarlo si heredan sus descriptores, normalmente después de un
\texttt{fork()}. Una FIFO (\emph{First In, First Out}), en cambio, tiene un
nombre y permite que procesos no relacionados la abran mediante ese nombre.

\begin{notebox}
Un pipe no es un archivo temporal. El núcleo mantiene un búfer para el canal y
entrega los bytes en el mismo orden en que fueron escritos. Si no queda ningún
descriptor de escritura abierto, una lectura puede observar fin de archivo;
si el búfer está vacío pero todavía existe un escritor, la lectura espera.
\end{notebox}

\subsection{Construir \texttt{ls | wc}}

El siguiente programa crea dos hijos. El primero ejecuta \texttt{ls} y redirige
su salida al pipe; el segundo ejecuta \texttt{wc} y redirige su entrada desde
el pipe. El proceso original cierra sus copias de ambos extremos y espera a
los dos hijos.

\clearpage
\enlargethispage{3pt}
\inputcode[language=C,basicstyle=\ttfamily\scriptsize]{../examples/tp5/tuberia.c}

Copiá el listado anterior en un archivo llamado \texttt{tuberia.c}. Luego
compilá y ejecutá el programa. La salida puede variar según los archivos del
directorio actual:

\begin{lstlisting}[language=bash]
$ gcc -Wall -Wextra -std=c11 -o tuberia tuberia.c
$ ./tuberia
      8      8     74
\end{lstlisting}

Los números mostrados son sólo un ejemplo. Probá también la versión del shell
y compará las dos ejecuciones:

\begin{lstlisting}[language=bash]
$ ls | wc
$ strace -f -e trace=pipe,pipe2,dup2,close,execve,wait4 ./tuberia
\end{lstlisting}

En el hijo que ejecuta \texttt{ls}, \texttt{dup2(fd[1], STDOUT\_FILENO)} hace
que el descriptor 1 apunte al extremo de escritura. En el otro hijo, la misma
operación hace que el descriptor 0 apunte al extremo de lectura. Después de
duplicar, cada proceso cierra los descriptores que ya no necesita. Esos
\texttt{close()} no son un detalle decorativo: un descriptor de escritura
olvidado puede impedir que \texttt{wc} vea el fin de archivo.

\begin{enumerate}
  \item Dibujá el pipe y marcá qué proceso conserva cada extremo después de
  los cierres. Indicá qué representan los descriptores 0 y 1 en cada hijo.
  \item Explicá por qué el programa crea dos hijos en vez de ejecutar
  \texttt{wc} directamente en el proceso original.
  \item Agregá una tercera etapa para construir una cadena equivalente a
  \texttt{ls | wc -l | cat}. ¿Qué extremos debe cerrar cada proceso?
  \item Consultá \texttt{/proc/<pid>/fd} mientras un proceso permanece
  bloqueado en una lectura. Relacioná los enlaces simbólicos con los
  descriptores del programa.
\end{enumerate}

\subsection{Tuberías con nombre: \texttt{mkfifo}}

Una FIFO tiene una entrada permanente en un directorio y puede ser abierta por
procesos que no tienen un ancestro común. El nombre no contiene los datos: sólo
permite que los procesos encuentren el mismo canal del núcleo. Abrirla para
leer o escribir puede bloquearse hasta que exista un proceso en el extremo
opuesto.

Probala desde dos terminales. En la primera terminal creá la FIFO y dejá un
lector esperando:

\begin{lstlisting}[language=bash]
$ mkfifo canal
$ ls -l canal
$ cat < canal
\end{lstlisting}

En la segunda terminal escribí un mensaje y observá que \texttt{cat} termina
cuando recibe el fin de archivo:

\begin{lstlisting}[language=bash]
$ printf 'mensaje enviado por una FIFO\n' > canal
\end{lstlisting}

La entrada \texttt{canal} permanece en el directorio aunque los procesos hayan
terminado. Eliminá la FIFO cuando finalices la prueba:

\begin{lstlisting}[language=bash]
$ rm canal
\end{lstlisting}

Repetí la prueba con \texttt{echo}, \texttt{wc} o dos programas propios.
Compará la FIFO con el pipe anónimo del programa anterior: ¿qué propiedad
permite que la usen procesos que no tienen un ancestro común?

\begin{extrabox}
Como extensión, inspeccioná los descriptores de los procesos mientras la FIFO
está abierta con \texttt{ls -l /proc/\textless pid\textgreater/fd}. Relacioná
los enlaces simbólicos con el nombre \texttt{canal} y con el estado bloqueado de
la operación de apertura.
\end{extrabox}

\section{Exclusión mutua con candados de archivos}

Cuando varios procesos acceden al mismo recurso, una operación puede necesitar
exclusión mutua: mientras un proceso está en su sección crítica, los demás
deben esperar. Sin coordinación, dos procesos pueden leer el mismo estado,
calcular resultados incompatibles y sobrescribir los cambios del otro.

\texttt{lockf()} permite solicitar un candado exclusivo sobre una región de
un archivo abierto. El tamaño se cuenta desde la posición actual del
descriptor. Las operaciones principales son:

\begin{center}
\begin{tabular}{c|l}
\textbf{Operación} & \textbf{Comportamiento} \\
\hline
\texttt{F\_LOCK} & solicita el candado y espera si la región está ocupada \\
\texttt{F\_TLOCK} & solicita el candado y falla inmediatamente si está ocupado \\
\texttt{F\_TEST} & comprueba si la región está disponible, sin tomarla \\
\texttt{F\_ULOCK} & libera el candado de la región
\end{tabular}
\end{center}

Estos candados son normalmente \emph{advisory}: el núcleo coordina a los
procesos que también intentan adquirir el candado, pero no impide que un
programa que lo ignora escriba el archivo. El candado se libera al ejecutar
\texttt{F\_ULOCK} o cuando el proceso cierra el archivo o termina.

\subsection{Un candado bloqueante}

El ejemplo toma el candado de los primeros 100 bytes durante cinco segundos.
El argumento opcional permite que varias copias usen el mismo archivo. La
llamada a \texttt{ftruncate()} deja preparada la región, y
\texttt{lseek()} coloca la posición al comienzo antes de llamar a
\texttt{lockf()}.

\inputcode[language=C,basicstyle=\ttfamily\footnotesize]{../examples/tp5/candado.c}

Copiá el listado anterior en un archivo llamado \texttt{candado.c}. Compilá el
archivo y ejecutá primero una sola copia:

\begin{lstlisting}[language=bash]
$ gcc -Wall -Wextra -std=c11 -o candado candado.c
$ ./candado archivo.bin
Proceso 2659: esperando el candado
Proceso 2659: candado colocado
Proceso 2659: candado eliminado
\end{lstlisting}

Después iniciá tres copias sobre el mismo archivo y esperá a que todas
terminen:

\begin{lstlisting}[language=bash]
$ ./candado archivo.bin & p1=$!
$ ./candado archivo.bin & p2=$!
$ ./candado archivo.bin & p3=$!
$ wait "$p1" "$p2" "$p3"
\end{lstlisting}

Sólo un proceso debería mostrar ``candado colocado'' a la vez. El orden de
adquisición no tiene por qué coincidir con el orden de lanzamiento ni con el
valor de los PID. Mientras el experimento está corriendo, observá
\texttt{/proc/locks} desde otra terminal:

\begin{lstlisting}[language=bash]
$ cat /proc/locks
\end{lstlisting}

Una línea puede indicar la clase del candado (por ejemplo, \texttt{POSIX}), si
es asesor o obligatorio, el modo de acceso, el PID, el dispositivo e inode y
los límites de la región. La representación y la cantidad de líneas dependen
de la versión del núcleo y de los demás procesos del sistema.

\begin{enumerate}
  \item Registrá para cada copia el PID, el instante en que empezó a esperar,
  el instante en que obtuvo el candado y el instante en que lo liberó.
  \item ¿Qué proceso entra primero? ¿Podés deducir una política de
  planificación a partir de ese orden? Justificá la respuesta.
  \item Cambiá \texttt{F\_LOCK} por \texttt{F\_TLOCK}. ¿Qué diferencia
  observás y qué código de error deberías tratar en un programa real?
  \item ¿Por qué este mecanismo no evita que un programa que use
  \texttt{write()} sin llamar a \texttt{lockf()} modifique el archivo?
\end{enumerate}

\begin{extrabox}
Investigá \texttt{flock()} y los candados POSIX de \texttt{fcntl()}. Compará
qué objetos identifican, qué ocurre al duplicar un descriptor y si los
candados son heredados después de \texttt{fork()}. Consultá las páginas de
manual de \texttt{flock(2)}, \texttt{fcntl(2)} y \texttt{lockf(3)}; no
modifiques archivos del sistema.
\end{extrabox}

\section{Semáforos de System V}

Un semáforo es un contador mantenido por el núcleo para coordinar el acceso a
un recurso. En el caso binario, el valor 1 puede representar ``disponible'' y
el valor 0 ``ocupado''. La operación de espera decrementa el contador y se
bloquea si no puede hacerlo; la operación de señal lo incrementa cuando el
proceso libera el recurso. A diferencia de un pipe, el semáforo no transporta
datos: sólo sincroniza.

En este apartado usaremos la interfaz de IPC (\emph{Inter-Process
Communication}) de System V:

\begin{itemize}
  \item \texttt{semget()} crea un conjunto de semáforos o devuelve el
  identificador de uno existente;
  \item \texttt{semop()} aplica una o más operaciones atómicas descritas por
  estructuras \texttt{sembuf};
  \item \texttt{semctl()} consulta, modifica o elimina el conjunto.
\end{itemize}

\subsection{Permisos del conjunto}

Los permisos de un semáforo de System V se aplican al conjunto completo y
siguen el esquema de propietario, grupo y otros usuarios. Sus nombres no son
``lectura'' y ``escritura'', sino lectura y alteración: no hay datos que leer o
escribir como en un archivo.

\begin{itemize}
  \item El permiso de lectura permite consultar valores y metadatos, por
  ejemplo con \texttt{GETVAL}, \texttt{GETALL} o \texttt{IPC\_STAT}.
  \item El permiso de alteración permite cambiar valores con \texttt{semop()},
  \texttt{SETVAL} o \texttt{SETALL}.
  \item \texttt{IPC\_SET} e \texttt{IPC\_RMID} requieren ser propietario o
  creador del conjunto, o tener privilegios suficientes.
\end{itemize}

Por ejemplo, \texttt{0660} permite leer y alterar el conjunto al propietario
y al grupo, pero no da acceso a otros usuarios. En la columna \texttt{perms}
de \texttt{ipcs -s}, el bit que en un archivo se suele llamar ``escritura''
representa aquí alteración.

Una estructura \texttt{sembuf} indica el índice del semáforo, el cambio que se
quiere hacer en \texttt{sem\_op} y opciones como \texttt{IPC\_NOWAIT}. Una
operación positiva incrementa el valor; una operación negativa espera hasta
que haya suficiente valor y luego lo decrementa. Una operación cero espera
hasta que el valor llegue a cero.

\subsection{Crear e inicializar un semáforo}

El programa siguiente solicita un conjunto con un semáforo. Usa la constante
\texttt{IPC\_PRIVATE} para pedir que se cree un conjunto nuevo. En este
contexto, esa constante no significa que el semáforo sea sólo del proceso:
indica a \texttt{semget()} que debe crear un conjunto. El programa incrementa el semáforo a 1 y deja el conjunto
existente para que puedas inspeccionarlo y eliminarlo en el paso siguiente.

\inputcode[language=C,basicstyle=\ttfamily\footnotesize]{../examples/tp5/creasem.c}

Copiá el listado anterior en un archivo llamado \texttt{creasem.c}. Compilá el
archivo y compará la lista de semáforos antes y después de ejecutar:

\begin{lstlisting}[language=bash]
$ gcc -Wall -Wextra -std=c11 -o creasem creasem.c
$ ipcs -s
$ ./creasem
Conjunto creado: 98307
Valor inicial: 1 (recurso disponible).
Eliminar con: ./cntrlsem 98307
$ ipcs -s
\end{lstlisting}

El identificador del ejemplo cambia en cada ejecución. Guardá el valor que
imprima tu sistema; no copies \texttt{98307}. Observá las columnas de
\texttt{ipcs}: clave, identificador, propietario, permisos y cantidad de
semáforos del conjunto.

\subsection{Consultar y eliminar el conjunto}

El comando \texttt{semctl()} recibe el identificador del conjunto y el índice
del semáforo. En este TP usamos un conjunto de un elemento, por eso el índice
es siempre cero. Algunas operaciones útiles son:

\begin{center}
\begin{tabular}{c|p{9.2cm}}
\textbf{Comando} & \textbf{Efecto} \\
\hline
\texttt{GETVAL} & devuelve el valor actual del semáforo \\
\texttt{SETVAL} & establece el valor indicado por el argumento \\
\texttt{GETPID} & devuelve el PID del último proceso que ejecutó \texttt{semop()} \\
\texttt{GETNCNT} & devuelve cuántos procesos esperan que el valor aumente \\
\texttt{GETZCNT} & devuelve cuántos procesos esperan que el valor sea cero \\
\texttt{GETALL} & devuelve los valores de todos los semáforos del conjunto \\
\texttt{SETALL} & establece los valores de todos los semáforos del conjunto \\
\texttt{IPC\_SET} & modifica propietario, grupo y permisos del conjunto \\
\texttt{IPC\_STAT} & copia la configuración del conjunto a una estructura \texttt{semid\_ds} \\
\texttt{IPC\_RMID} & elimina el conjunto de semáforos
\end{tabular}
\end{center}

\texttt{GETALL} y \texttt{SETALL} usan un arreglo dentro del argumento de
\texttt{semctl()}; \texttt{IPC\_STAT} y \texttt{IPC\_SET} usan una estructura
\texttt{semid\_ds}. En todos los casos, el proceso debe tener los permisos
necesarios para consultar o modificar el conjunto.

El programa de control usa \texttt{IPC\_RMID}:

\inputcode[language=C,basicstyle=\ttfamily\scriptsize]{../examples/tp5/cntrlsem.c}

Copiá el listado anterior en un archivo llamado \texttt{cntrlsem.c} y compilá el
archivo:

\begin{lstlisting}[language=bash]
$ gcc -Wall -Wextra -std=c11 -o cntrlsem cntrlsem.c
\end{lstlisting}

Pasá el identificador que imprimió \texttt{creasem} y verificá que ya no
aparezca en \texttt{ipcs -s}:

\begin{lstlisting}[language=bash]
$ ./cntrlsem 98307
Conjunto eliminado.
$ ipcs -s
\end{lstlisting}

Si cerraste la terminal antes de eliminar el conjunto, todavía podés quitarlo
con \texttt{ipcrm -s SEMID}, siempre que tengas permisos. Los semáforos de
System V son recursos persistentes del núcleo: terminar el programa creador no
los elimina automáticamente.

\begin{enumerate}
  \item Ejecutá \texttt{creasem} dos veces sin eliminar los conjuntos.
  ¿Cuántas entradas nuevas aparecen y por qué \texttt{IPC\_PRIVATE} no hace
  que ambas ejecuciones compartan el mismo conjunto?
  \item Eliminá los dos conjuntos y repetí la observación con
  \texttt{ipcs -s}. ¿Qué diferencia hay entre el valor del semáforo y el
  identificador del conjunto?
  \item Relacioná la operación negativa de \texttt{semop()} con la espera
  bloqueante de \texttt{lockf(F\_LOCK)}. ¿Qué tienen en común y qué no
  protegen por sí solos?
  \item Consultá \texttt{/proc/sysvipc/sem} y compará su información con
  \texttt{ipcs -s}.
\end{enumerate}

\end{document}
